\documentclass[11pt,a4paper]{article}
\usepackage[T1]{fontenc}
\usepackage{lmodern,amsmath,amssymb,amsthm,mathtools,booktabs,hyperref}
\usepackage[margin=27mm]{geometry}
\setlength{\emergencystretch}{2em}
\newtheorem{theorem}{Theorem}[section]
\newtheorem{proposition}[theorem]{Proposition}
\newtheorem{lemma}[theorem]{Lemma}
\newtheorem{corollary}[theorem]{Corollary}
\theoremstyle{definition}\newtheorem{definition}[theorem]{Definition}
\newcommand{\Q}{\mathbb Q}\newcommand{\Z}{\mathbb Z}
\newcommand{\F}{\mathbb F}
\title{No constant pole parameter is complete}
\author{}\date{}
\begin{document}\maketitle
\begin{abstract}
We consider the two-equation pole certificate used in the rational-source
construction of the integers.  No fixed rational value of its parameter
$e$ is complete on the prescribed set of prepared inputs.  For either
degenerate value $e=1,-1/3$, the same prepared input
$b=(204/55)^2$ gives a genus-one quartic with points over every completion
of $\Q$ and no rational point.  Exact unconditional two-descent and
Cassels-pairing computations certify the two elliptic rank bounds used
in the proof.  An elementary squareclass argument excludes every other
constant value of $e$.  The obstruction includes zero values of the
second square root, so removing a nonzero guard does not repair it.
\end{abstract}

\section{Prepared inputs and the pole equations}

\begin{definition}
A rational number $b$ is \emph{prepared} if its reduced denominator is
a square and
\[
 \begin{gathered}
 0<b<17,\qquad b\in\Q_2^{\times2}\cap\Q_3^{\times2}
                      \cap\Q_{17}^{\times2},\\
 v_2(b)\ge3,\qquad v_3(b)\ge1,\qquad v_{17}(b)\ge2.
 \end{gathered}
\]
For such $b$, put $B=17b$ and $V=B(B-b^2)$.  The pole equations are
\begin{equation}\label{eq:pole}
 p^2=B(1-e)(1+3e),\qquad
 Bq^2=(s^2-eV)(s^2+3eV).
\end{equation}
All four witnesses are rational; zero values are allowed.
\end{definition}

These are exactly the preparation hypotheses and the zero-preparation
equations stated in~\cite[Section~3]{Dependencies}.  Its realization
theorem permits $e$ to vary.  Our question is whether a fixed value
can replace that choice.

\begin{theorem}\label{thm:fixed-e}
For every $e_0\in\Q$, there is a prepared $b$ for which
\eqref{eq:pole} has no rational solution with $e=e_0$.
For $e_0=1$ and $e_0=-1/3$, one can take
\[
                          b_0=\frac{41616}{3025}=\left(\frac{204}{55}\right)^2.
\]
For either of these two values, the second equation has no rational
solution even though its smooth projective quartic has points over
every completion of $\Q$.
\end{theorem}

The rank bounds in the proof are exact computations.  They use the
unconditional algebraic rank bounds of PARI/GP, not an analytic-rank
calculation or the Birch--Swinnerton-Dyer conjecture.  The complete
verifier and its exact intermediate data accompany this note.

\section{Two elliptic rank certificates}

For $n\in\Q^\times$, write
\[
                 E_n:\quad y^2=x(x-n)(x+3n).
\]
These curves have the three rational nonidentity two-torsion points
$(0,0),(n,0),(-3n,0)$.

\begin{proposition}[Exact computation]\label{prop:ranks}
The groups $E_{-29427}(\Q)$ and $E_{9809}(\Q)$ have ranks zero
and one, respectively.  Both have torsion subgroup
$(\Z/2\Z)^2$.  The point
\[
                      P_*=(-20449,2356926)\in E_{9809}(\Q)
\]
has infinite order.
\end{proposition}
\begin{proof}
The verifier initializes the two curves with coefficient vectors
$[0,2n,0,-3n^2,0]$.  PARI/GP version $2.17.2$ returns
\[
 \begin{array}{c|c}
 n&\texttt{ellrank}\text{ output}\\\hline
 -29427&[0,0,2,[]]\\
 9809&[1,1,2,[[-20449,2356926]]].
 \end{array}
\]
Here the first two entries are unconditional lower and upper bounds
for the rational rank.  The algorithm computes the two-Selmer rank,
the rational two-torsion rank, and the rank of the relevant Cassels
pairing; the upper bound is the first rank minus the other two
\cite{PARI}.  The saved exact diagnostics give two-Selmer ranks $4,5$
and pairing matrices
\[
 C_-=
 \begin{pmatrix}
 0&1&0&1\\1&0&1&0\\0&1&0&1\\1&0&1&0
 \end{pmatrix},\qquad
 C_+=
 \begin{pmatrix}
 0&0&0&0&0\\
 0&0&1&0&1\\
 0&1&0&1&1\\
 0&0&1&0&1\\
 0&1&1&1&0
 \end{pmatrix}.
\]
Independent row reduction over $\F_2$ verifies that each has rank
two.  Since the rational two-torsion has rank two, the certified
upper bounds are $4-2-2=0$ and $5-2-2=1$.

For an independent torsion check, direct finite-field counting gives
\[
 \begin{array}{c|rrrr}
 n&\#E_n(\F_5)&\#E_n(\F_7)&\#E_n(\F_{11})&\#E_n(\F_{13})\\\hline
 -29427&4&8&8&12\\
 9809&8&8&16&12.
 \end{array}
\]
All these primes are good.  The injection on prime-to-$p$ torsion
under good reduction, using the different primes in the table,
bounds the rational torsion order by four in each case.  The known
two-torsion attains this bound.  Direct substitution verifies $P_*$,
and its nonzero ordinate shows that it is not torsion.  This also
independently supplies the lower rank bound one for $E_{9809}$.
\end{proof}

\begin{lemma}\label{lem:x-class}
For a nonsingular curve $y^2=x(x^2+ax+c)$ with $c\ne0$, the map
\[
 \alpha:E(\Q)\longrightarrow\Q^\times/\Q^{\times2},\qquad
 \alpha(O)=1,\quad\alpha((0,0))=c,\quad
 \alpha((x,y))=x\quad(x\ne0)
\]
is a group homomorphism.  It kills $2E(\Q)$.  In particular, if
$E$ has rank $r$ and full rational two-torsion, its image has size
at most $2^{r+2}$.
\end{lemma}
\begin{proof}
Intersect a nonvertical line $y=ux+v$ with the curve.  The product
of its three intersection abscissas, counted with multiplicity,
is $v^2$.  Since reflection in the $x$-axis does not change an
abscissa, this proves the homomorphism identity when $v\ne0$.
If $v=0$, the two nonzero intersection abscissas have product $c$,
which gives the identity when the third point is $(0,0)$.
Addition by $(0,0)$ sends a nonzero abscissa $x$ to $c/x$;
vertical lines and two-torsion doubling give the remaining cases.
The target has exponent two, so $2E(\Q)$ is killed.  Finally the
Mordell--Weil decomposition gives
$|E(\Q)/2E(\Q)|=2^r|E(\Q)[2]|=2^{r+2}$.
\end{proof}

\begin{corollary}\label{cor:class17}
Neither $E_{-29427}(\Q)$ nor $E_{9809}(\Q)$ has a nonzero
abscissa in the rational squareclass of $17$.
\end{corollary}
\begin{proof}
The first curve has only its four torsion points.  Its nonzero
abscissas are $-29427$ and $88281=9\cdot17\cdot577$; neither has
squareclass $17$.

On $E_{9809}$, the two-torsion points give squareclasses $-3$ and
$9809$, while $P_*$ gives $-1$, because $20449=143^2$.
These three classes are independent: $9809=17\cdot577$ and the
prime $577$ distinguishes the third generator from the first two.
They therefore supply eight classes.  By Lemma~\ref{lem:x-class}
and rank one, these exhaust the image:
\[
          \alpha(E_{9809}(\Q))
          =\{\pm1,\ \pm3,\ \pm9809,\ \pm29427\}.
\]
The class $17$ is absent.
\end{proof}

\section{The two degenerate specializations}

Put $M=55$, $t=204/M$, and $d=M^2-2448=577$.  Then
\[
 b_0=t^2,\qquad B=17t^2,\qquad
 17-t^2=\frac{17d}{M^2},\qquad V=\frac{289t^4d}{M^2}.
\]
For $e=1$, the invertible rational substitutions
\[
             s=\frac{17t^2}{M}u,\qquad
             q=\frac{17t^3}{M^2}v
\]
turn the second pole equation into
\begin{equation}\label{eq:Cplus}
                  C_+:\quad v^2=17(u^2-d)(u^2+3d).
\end{equation}
For $e=-1/3$, the substitutions
\[
             s=\frac{17t^2}{3M}u,\qquad
             q=\frac{17t^3}{9M^2}v
\]
give
\begin{equation}\label{eq:Cminus}
                  C_-:\quad v^2=17(u^2+3d)(u^2-9d).
\end{equation}
The verifier checks these identities coefficient by coefficient.

In either equation $u=0$ makes the right side negative.  Every
rational affine point would therefore have $u\ne0$.  The map
\[
                         (x,y)=(17u^2,17uv)
\]
sends $C_+$ to $E_{17d}=E_{9809}$ and $C_-$ to
$E_{-51d}=E_{-29427}$.  Its nonzero abscissa has squareclass $17$,
contradicting Corollary~\ref{cor:class17}.  This proves that both
quartics have no rational affine point.  Their points at infinity
would require $17$ to be a rational square, so their smooth
projective completions have no rational point either.

\begin{proposition}\label{prop:local}
Both $C_+$ and $C_-$ have points over every completion of $\Q$.
\end{proposition}
\begin{proof}
Let $g_\pm(u)$ denote the corresponding right side.  The following
table gives rational values of $u$ for which $g_\pm(u)$ is a
nonzero square in the indicated local field.  Each entry records
$u$, the valuation, and the residue of the unit after removing
that valuation.  The unit is taken modulo $8$ at $2$ and modulo
$p$ at an odd prime.
\[
\begin{array}{c|c|c}
 p&C_+:(u,v_p(g_+),\text{unit})&C_-:(u,v_p(g_-),\text{unit})\\\hline
 2&(23,6,1)&(21,6,1)\\
 3&(1,2,1)&(12,4,1)\\
 17&(4,2,13)&(22,2,9)\\
 577&(1,0,17)&(1,0,17).
\end{array}
\]
The units $13,9$ are squares modulo $17$, and $17$ is a square
modulo $577$.  The standard local square criterion proves these
four local assertions.  At the real place, $u=100$ makes both
right sides positive.

The quartic discriminants are
\[
 \operatorname{disc}(g_+)=-2^{12}3\,17^6d^6,
 \qquad
 \operatorname{disc}(g_-)=-2^{12}3^7\,17^6d^6.
\]
Thus the smooth projective genus-one curves have good reduction
outside $2,3,17,577$.  At any other prime, the Hasse bound gives
a point on the reduction, and smooth lifting gives a local point.
This proves the proposition.
\end{proof}

\section{Exclusion of all other constant values}

The number $b_0=(204/55)^2$ is prepared: its reduced denominator
is square, it lies in $(0,17)$, and its valuations at $2,3,17$
are $4,2,2$.  Being a rational square, it satisfies all three
local square conditions.

A second prepared number is
\[
 b_1=409\left(\frac{204}{1001}\right)^2
       =\frac{17020944}{1002001}.
\]
It is less than $17$ because $17020944<17034017$.
Its denominator is square and its valuations at $2,3,17$ are
again $4,2,2$.  The number $409$ is $1$ modulo $8,3,17$,
so it is a square in all three local fields.

Suppose $e_0\ne1,-1/3$ made the first pole equation soluble
for both inputs.  Its factor $(1-e_0)(1+3e_0)$ is nonzero;
division of the two equations would make $b_1/b_0$ a rational
square.  But
\[
                     \frac{b_1}{b_0}=409\left(\frac{55}{1001}\right)^2
\]
is not a square.  Thus at least one of the two prepared inputs
already fails the first equation.  The two excluded degenerate
values fail at $b_0$ by Section~3.  This completes the proof of
Theorem~\ref{thm:fixed-e}.

The theorem concerns a constant $e$.  It does not exclude a
nonconstant function of $b$, a different relation between the pole
parameters, or a different pole predicate.  Those proposals require
their own completeness arguments.

\section*{Reproducible verification}

The program
\begin{center}\small
\nolinkurl{check_fixed_e_obstruction.py}
\end{center}
generates \nolinkurl{fixed-e-obstruction-certificate.json} and
\nolinkurl{fixed-e-pari-descent.log}.  Apart from the two algebraic
rank computations, it uses exact integer arithmetic and rational
fractions from the Python standard library.  The preparation
conditions, substitutions, torsion bounds, squareclass exhaustion,
local witnesses, and ranks of the saved pairing matrices are checked
independently of the rank routine.  A rational-point search finding
no points is not used as a nonexistence certificate.

\begin{thebibliography}{9}
\bibitem{Dependencies}
\emph{Arithmetic dependencies of the construction with 29 variables},
project paper, version 8, 8 November 2026.
\bibitem{PARI}
The PARI Group, \emph{PARI/GP: Elliptic curves}, documentation of
\texttt{ellrank}, including its unconditional two-descent and
Cassels-pairing bounds,
\url{https://pari.math.u-bordeaux.fr/dochtml/html/Elliptic_curves.html\#ellrank}.
\end{thebibliography}
\end{document}
